
% ============================================================
%  REMERCIEMENTS
% ============================================================
\clearpage
\thispagestyle{empty}
\chapter*{Remerciements}

Au terme de ce travail, je ressens une profonde reconnaissance envers toutes les
personnes qui ont contribué, de près ou de loin, à l'aboutissement de ce projet.
Un mémoire de fin d'études n'est jamais uniquement le fruit d'un effort
individuel : il porte aussi les traces des conseils reçus, des encouragements
silencieux, des sacrifices familiaux et des moments de doute surmontés grâce à
la présence des autres.

Je tiens d'abord à exprimer ma gratitude à mon encadrant académique,
\textbf{Mr Nadjib Bendaoud}, pour son accompagnement, ses orientations et la
confiance qu'il m'a accordée tout au long de ce parcours. Ses remarques et ses
conseils ont contribué à structurer ma démarche et à donner à ce travail une
forme plus rigoureuse et plus cohérente.

J'adresse également mes sincères remerciements à mon encadrante professionnelle,
\textbf{M jihen hasnaoui}, ainsi qu'à toute l'équipe d'\textbf{Intech Solutions},
pour leur accueil, leur disponibilité et l'environnement professionnel dans
lequel ce projet a pu évoluer. Leur soutien m'a permis de confronter mes idées
aux réalités du terrain et de transformer une proposition technique en une
solution concrète, utile et adaptée au contexte tunisien.

Je remercie chaleureusement l'ensemble de mes enseignants à \textbf{TEK-UP}, qui
ont participé à ma formation et m'ont transmis les connaissances nécessaires
pour aborder ce projet avec méthode, exigence et ambition. Chaque cours, chaque
remarque et chaque défi rencontré durant ces années ont contribué à construire
la personne et l'ingénieur que je deviens aujourd'hui.

Mes remerciements les plus profonds vont à ma famille. À mes parents, en
particulier, je dois bien plus que des mots ne peuvent exprimer. Leur patience,
leurs sacrifices, leurs prières et leur confiance ont été mon refuge dans les
moments de fatigue et ma force dans les périodes d'incertitude. Si ce travail
marque une étape importante de mon parcours, il est aussi le reflet de tout ce
qu'ils ont donné sans jamais rien attendre en retour.

Je remercie enfin mes amis, mes collègues et toutes les personnes qui m'ont
soutenu durant cette période exigeante. Une parole simple, un encouragement au
bon moment ou une présence sincère peuvent parfois suffire à rallumer la
motivation lorsque le chemin semble trop long.

À toutes et à tous, je dédie ma reconnaissance la plus sincère. Ce travail est
le résultat d'un parcours fait d'efforts, d'épreuves et d'espoir ; il n'aurait
pas eu la même valeur sans votre soutien.
